\section{Hinduism}

Hinduism is the family of traditions rooted in the Vedas and developed through the Upanishads, the epics, the Puranas, and the later schools of Vedanta and Tantra. It is not a single doctrine but a set of related systems that share a common vocabulary and disagree on which reality is ultimate. That internal disagreement makes it especially useful for comparison, since it holds several rival models of the absolute, the self, and the world in one tradition.

\subsection{The Creation Model}

There is no single creation story. The oldest songs give several accounts, and the tradition lets them stand together without forcing them to agree. One common shape appears across them.

\subsection{In the beginning, neither being nor non-being}

Before anything, there was neither being nor non-being. There was no sky, no air, no death, no deathlessness, no day, no night. There was only \textbf{the One}, which breathed without breath, by its own power, depending on nothing outside itself. Beyond it there was nothing. Darkness lay hidden within darkness. Everything was an undivided water with no shore and no surface.

Then something stirred. The hymns name two forces. First a \textbf{heat}, the pressure that drives one thing to separate from another. Then a \textbf{desire}, the first motion of mind reaching toward something other than itself. This desire was the first seed. With it, the undivided began to incline toward division. The thinkers who searched their own hearts reported finding there the cord that ties what exists to what does not.

\subsection{The one becomes the many}

Once the stirring began, the One started to become countable. The whole became the world.

One account describes this as a \textbf{great offering}. There was a vast cosmic person, with countless heads and eyes and feet, who covered the whole earth and reached past it, who was everything that had been and everything still to come. This person was given up and divided, and from the pieces the world was assembled. From the mind came the moon. From the eye came the sun. From the breath came the wind. From the head came the sky and from the feet the earth. Even the ranks of human society came out of the body. Creation in this account is the whole broken into parts, the one divided into the many.

Another account describes a \textbf{golden womb} or golden egg. It arose first of all, floating on the formless water, holding within it the seed of everything that would be born. It became the single lord of all that came to be. It held up the earth and the sky and gave out life and breath. Every verse of this hymn returns to the same question, asking which power deserves the offering. The hymn is a search for the one behind the many.

In a later account the egg cracks open into the heavens and the earth, and from it a \textbf{maker} is born who orders the worlds, the beings, and the first ancestors. This maker does not build from nothing. He re-forms and re-orders what had dissolved at the close of the previous round.

\subsection{The world breathed out and breathed in}

Creation is not a single event. It is one phase of a continuing rhythm. The worlds are breathed out, held for an immense span, then breathed back in. Three movements repeat. A \textbf{bringing forth} at the dawn of each cosmic day. A long \textbf{holding} that keeps the worlds in order. And a \textbf{dissolving} at the close of each cosmic night, when everything returns to the unmanifest. After the night, the worlds form again. The spans involved run to billions and then trillions of years, and at the end of the longest cycle all of it returns fully to the source, until a new round begins.

\subsection{The honesty at the end}

The oldest hymn ends without claiming certainty. After describing the beginning, it asks who could really know where creation came from. The ordering powers, the gods themselves, came after the world was already made, so even they were not present to see it. The hymn closes by saying that perhaps the One who watches from the highest place knows how it began, or perhaps even he does not know. The account is explicit that the first moment remains hidden.
